\documentclass[11pt]{article}
\usepackage[margin=1in]{geometry}
\usepackage{amsmath,amssymb,booktabs,array,hyperref,microtype,xcolor}
\usepackage{enumitem}
\setlist[itemize]{leftmargin=*,topsep=2pt,itemsep=2pt}

\title{PNP Log-Concentration Transition\\
       \large Side-by-Side Comparison with the Concentration Formulation}
\author{Jake Weinstein}
\date{Week of April~27, 2026}

\newcommand{\Eeq}{E_{\mathrm{eq}}}
\newcommand{\VT}{V_{T}}
\newcommand{\Gelec}{\Gamma_{\mathrm{elec}}}
\newcommand{\Gconc}{\Gamma_{\mathrm{conc}}}
\newcommand{\Gground}{\Gamma_{\mathrm{ground}}}
\newcommand{\dx}{\,\mathrm{d}x}
\newcommand{\ds}{\,\mathrm{d}s}

\begin{document}
\maketitle

\section{Purpose and scope}
The active production stack
(\verb|Forward/bv_solver/forms_logc.py|) uses
$u_i = \ln c_i$ as primary unknowns instead of $c_i$. This note maps every
equation in \texttt{docs/PNP Equation Formulations.tex} to its
log-concentration counterpart so the transition can be verified line by line.
The mathematical change is a single substitution; the practical change is the
elimination of negative-concentration Newton iterates and the resulting
floor/regularization terms. The Poisson equation is structurally unchanged
modulo $c_i \to e^{u_i}$ in its source.

Throughout, the substitution
\[
  c_i = e^{u_i},
  \qquad
  \nabla c_i = e^{u_i}\,\nabla u_i = c_i\,\nabla u_i,
\]
is applied to every term in the strong and weak forms. The combination
$\nabla c_i + \tfrac{z_i F}{RT}\,c_i\,\nabla\phi$ that appears in the
Nernst--Planck flux factors as $c_i \bigl(\nabla u_i + \tfrac{z_i
F}{RT}\,\nabla\phi\bigr)$, which is the central simplification.

\section{Strong form}

\begin{center}
\renewcommand{\arraystretch}{1.6}
\begin{tabular}{@{} >{\bfseries}p{0.18\textwidth} p{0.36\textwidth} p{0.40\textwidth} @{}}
\toprule
& \textbf{Concentration ($c_i$)} & \textbf{Log-concentration ($u_i = \ln c_i$)} \\
\midrule
NP equation &
$\displaystyle \frac{\partial c_i}{\partial t}
   = \nabla\!\cdot\!\Bigl(D_i\bigl[\nabla c_i
       + \tfrac{z_i F}{RT}\,c_i\,\nabla\phi\bigr]\Bigr)$ &
$\displaystyle \frac{\partial e^{u_i}}{\partial t}
   = \nabla\!\cdot\!\Bigl(D_i\,e^{u_i}\bigl[\nabla u_i
       + \tfrac{z_i F}{RT}\,\nabla\phi\bigr]\Bigr)$ \\
NP flux &
$\mathbf J_i = -D_i\bigl(\nabla c_i + \tfrac{z_i F}{RT}\,c_i\,\nabla\phi\bigr)$ &
$\mathbf J_i = -D_i\,e^{u_i}\bigl(\nabla u_i + \tfrac{z_i F}{RT}\,\nabla\phi\bigr)$ \\
Poisson &
$-\nabla\!\cdot\!(\epsilon\,\nabla\phi) = \sum_i z_i F\,c_i$ &
$-\nabla\!\cdot\!(\epsilon\,\nabla\phi) = \sum_i z_i F\,e^{u_i}$ \\
BV flux BC on $\Gelec$ &
$-D_i\bigl(\nabla c_i + \tfrac{z_i F}{RT}\,c_i\,\nabla\phi\bigr)\cdot\mathbf n = \pm \tfrac{j}{n_e F}$ &
$-D_i\,e^{u_i}\bigl(\nabla u_i + \tfrac{z_i F}{RT}\,\nabla\phi\bigr)\cdot\mathbf n = \pm \tfrac{j}{n_e F}$ \\
Bulk concentration BC on $\Gconc$ &
$c_i = c_{0,i}$ &
$u_i = \ln c_{0,i}$ \quad (with floor $\max(c_{0,i},10^{-20})$ for product species with $c_{0,i}=0$) \\
\bottomrule
\end{tabular}
\end{center}

\noindent
\textbf{What is unchanged}: the Poisson PDE structure, the dielectric
$\epsilon$, the valences $z_i$, $E_{\mathrm{eq}}$, the bulk Poisson source
sum, and the choice of which boundary carries the BV flux. Only the
representation of $c_i$ has changed.

\textbf{What is new}: the Dirichlet bulk BC is now $u_i = \ln c_{0,i}$,
which is undefined for $c_{0,i}=0$. The code uses a floor
$c_{0,i} \leftarrow \max(c_{0,i}, 10^{-20})$ for product species
(\texttt{forms\_logc.py:441,510}); this matters for $\mathrm{H_2O_2}$,
whose physical bulk IC is zero.

\section{Weak form}

Multiplying the NP equation by $v\in V$ and integrating by parts (with the
electrode boundary handled by the BV flux condition):

\begin{center}
\renewcommand{\arraystretch}{1.5}
\begin{tabular}{@{} >{\bfseries}p{0.13\textwidth} p{0.40\textwidth} p{0.42\textwidth} @{}}
\toprule
& \textbf{Concentration} & \textbf{Log-concentration} \\
\midrule
NP weak form &
$\displaystyle\int_\Omega \frac{\partial c_i}{\partial t}\,v\dx
 + \int_\Omega D_i\bigl[\nabla c_i + \tfrac{z_i F}{RT}\,c_i\,\nabla\phi\bigr]\!\cdot\!\nabla v\dx$
\newline\hspace*{4em}$=\displaystyle\int_{\Gelec}\!\Bigl(\!\mp\tfrac{j}{n_e F}\Bigr) v\ds$ &
$\displaystyle\int_\Omega \frac{\partial e^{u_i}}{\partial t}\,v\dx
 + \int_\Omega D_i\,e^{u_i}\bigl[\nabla u_i + \tfrac{z_i F}{RT}\,\nabla\phi\bigr]\!\cdot\!\nabla v\dx$
\newline\hspace*{4em}$=\displaystyle\int_{\Gelec}\!\Bigl(\!\mp\tfrac{j}{n_e F}\Bigr) v\ds$ \\
\midrule
Poisson weak form &
$\displaystyle\int_\Omega \epsilon\,\nabla\phi\!\cdot\!\nabla w\dx
 = \int_\Omega\!\Bigl(\sum_i z_i F\,c_i\Bigr) w\dx$ &
$\displaystyle\int_\Omega \epsilon\,\nabla\phi\!\cdot\!\nabla w\dx
 = \int_\Omega\!\Bigl(\sum_i z_i F\,e^{u_i}\Bigr) w\dx$ \\
\bottomrule
\end{tabular}
\end{center}

\subsection{Time discretization (backward Euler)}
The code uses backward Euler in $c$-space, even though the unknown is $u$:
\[
  \boxed{\;
  \int_\Omega \frac{e^{u_i^{n+1}} - e^{u_i^{n}}}{\Delta t}\,v\dx
  \;\neq\;
  \int_\Omega e^{u_i^{n+1}}\frac{u_i^{n+1} - u_i^{n}}{\Delta t}\,v\dx
  \;}
\]
The left-hand form (used in \verb|forms_logc.py:294|) preserves discrete
species mass directly; the right-hand form would not. This is the only point
in the discretization where the choice of variable could in principle drift
from the concentration formulation, and the code is correct.

\section{Butler--Volmer boundary residual}

The BV reaction rate is
\[
  R_{\mathrm{BV}}
  = k_0\Bigl[c_{\mathrm{cat}}\,e^{-\alpha n_e \eta/V_T}
            - c_{\mathrm{anod}}\,e^{(1-\alpha)n_e \eta/V_T}\Bigr],
  \qquad
  \eta = \phi_m - \phi_s - \Eeq.
\]
Both branches in \verb|forms_logc.py| evaluate the same rate
mathematically; they differ in how they assemble it numerically.

\subsection{Standard branch (\texttt{bv\_log\_rate=False})}
Evaluates surface concentrations directly from the (clamped) log unknowns:
\[
  c_{\mathrm{surf},i} = \exp\!\bigl(\mathrm{clamp}(u_i,\,\pm 30)\bigr),
  \qquad
  R = k_0\,c_{\mathrm{surf,cat}}\,e^{-\alpha n_e\eta/V_T}
        - k_0\,c_{\mathrm{surf,anod}}\,e^{(1-\alpha)n_e\eta/V_T}.
\]
This is the literal port of the original BV formulation to the log-c
unknowns. The clamp $\pm 30$ is required because $u$ can wander wildly
during Newton iteration and \texttt{exp(u)} would otherwise overflow.

\subsection{Log-rate branch (\texttt{bv\_log\_rate=True}, current default)}
Algebraically rearranges the cathodic and anodic terms so the only $\exp$
applies to a single scalar sum, deferring evaluation:
\begin{align*}
  \log r_{\mathrm{cath}} &= \ln k_0 + u_{\mathrm{cat}}
       + \sum_{j} p_j\bigl(u_j - \ln c_{\mathrm{ref},j}\bigr)
       - \alpha\,n_e\,\eta/V_T \\
  \log r_{\mathrm{anod}} &= \ln k_0 + u_{\mathrm{anod}}
       + (1-\alpha)\,n_e\,\eta/V_T \\
  R &= e^{\log r_{\mathrm{cath}}} - e^{\log r_{\mathrm{anod}}}.
\end{align*}
The $u_i$ enter \emph{unclamped}. There is no $c_{\mathrm{surf}}$ floor in
this branch. This is what eliminates the phantom $\mathrm{R}_2$ sink at
$V_{\mathrm{RHE}} > +0.30$ V (when $u_{\mathrm{H_2O_2}}$ approaches the
$-30$ floor of the standard branch).

The two branches are identical in exact arithmetic; the difference is
purely numerical. See \texttt{docs/clipping\_conventions.md} for the full
discussion.

\section{Active clamps in the production solver}

The transition to log-c does \emph{not} eliminate clipping. Three distinct
guards exist in the code; only one was removed by the log-rate algebraic
rearrangement.

\begin{center}
\renewcommand{\arraystretch}{1.4}
\small
\begin{tabular}{@{} p{0.27\textwidth} p{0.27\textwidth} c p{0.30\textwidth} @{}}
\toprule
\textbf{Clip / clamp} & \textbf{Where in code} & \textbf{Active?} & \textbf{Guards against} \\
\midrule
$\eta$-clip $\pm 50$ on $\eta_{\mathrm{scaled}}$ &
\verb|_build_eta_clipped| (l.226) & yes &
overflow of $\exp(\alpha n_e\eta)$ \\
$u$-clamp $\pm 30$ on bulk $u_i$ &
bulk PDE residual (l.203) & yes &
overflow / underflow of $e^{u_i}$ as PDE coefficient \\
$c_{\mathrm{surf}}$ clamp inside BV &
legacy BV branch only (l.369--389) & no (log-rate) &
under-floored $e^{u}$ corrupting BV at high $V$ \\
\bottomrule
\end{tabular}
\end{center}

\noindent
\textbf{Why the bulk $u$-clamp cannot be eliminated by the log-rate trick}:
in the BV residual the entire reaction rate collapses to a single
$\exp(\text{scalar sum})$, so the $\exp$ can be deferred. In the bulk PDE,
$e^{u_i}$ appears as a \emph{coefficient} multiplying $\nabla v$ and other
terms that vary across the domain --- there is no single $\exp$ to defer.
The $u$-clamp is therefore floating-point safety only; widen to $u_{\max}=100$
when $V_{\mathrm{RHE}} > +0.30$ V because $u_{\mathrm{H_2O_2}}$ at SS
itself approaches $-30$ in that regime.

\noindent
\textbf{Why the $\eta$-clip cannot be eliminated either}:
$\exp(\pm 138)\sim 10^{\pm 60}$ overflows double precision; $\eta$-clip at
$\pm 50$ is the standard guard. R$_2$ ($\Eeq = 1.78$~V, $n_e = 2$) unclips
when $|V - \Eeq| < 50\,V_T = 1.285$~V, i.e. $V > +0.495$~V. R$_1$
($\Eeq = 0.68$~V, $n_e = 1$) unclips for $V > -0.605$~V, so it is
essentially always unclipped on the production grid.

\section{Verification checklist}
Items to spot-check while reading the code against this writeup. All line
numbers refer to \texttt{Forward/bv\_solver/forms\_logc.py} unless otherwise
noted.

\begin{enumerate}\setlength\itemsep{4pt}
  \item \textbf{Flux factoring} (l.\,291). \\
        \verb|Jflux = D[i] * c_i * (grad(u_i) + grad(drift))| \\
        Confirms $\mathbf J = -D\,c\,(\nabla u + \tfrac{zF}{RT}\nabla\phi)$.
  \item \textbf{Time derivative} (l.\,294). \\
        \verb|((c_i - c_old) / dt_const) * v * dx| \\
        Confirms $\Delta(e^u)/\Delta t$, not $e^u\,\Delta u/\Delta t$.
  \item \textbf{Poisson source} (l.\,428). \\
        \verb|sum(z[i] * ci[i] * w for i in range(n))| \\
        Confirms $\sum_i z_i\,e^{u_i}$; the $F$ prefactor is folded
        into \verb|charge_rhs|.
  \item \textbf{Bulk Dirichlet BC} (l.\,445--447). \\
        \verb|c0_i = max(float(c0_model[i]), _C_FLOOR)| then
        \verb|fd.Constant(np.log(c0_i))|. \\
        Confirms $u_i = \ln c_{0,i}$ on $\Gconc$ with floor
        $c_{0,i}\leftarrow\max(c_{0,i}, 10^{-20})$ for product species.
  \item \textbf{Initial conditions} (l.\,511--513). Same floor convention as
        the bulk BC; assigns $u_i^{(0)} = \ln c_{0,i}$.
  \item \textbf{No blob IC.} The kwarg
        \verb|set_initial_conditions(..., blob=True)| is silently ignored
        in log-c (see \verb|dispatch.py:32-41|). The blob IC was a
        concentration-formulation feature with no log-c analog.
  \item \textbf{$\eta$-clip on $\eta_{\mathrm{scaled}}$, not on
        $\alpha n_e \eta_{\mathrm{scaled}}$} (l.\,244). \\
        \verb|min_value(max_value(eta_scaled, -clip_val), clip_val)| \\
        Applied \emph{before} multiplying by $\alpha n_e$. This is the
        source of the $V_{\mathrm{RHE}} > +0.495$~V R$_2$ unclip
        threshold (\emph{not} the $+1.14$~V number from older handoffs).
  \item \textbf{Log-rate BV uses unclamped $u_i$} (l.\,337). The cathodic
        log-rate is built as
        \begin{quote}\small
          \verb|log_cathodic = fd.ln(k0_j) + ui[cat_idx]|\\
          \hspace*{2em}\verb|- alpha_j * n_e_j * eta_j|
        \end{quote}
        Compare to the legacy branch (l.\,369) which uses
        \verb|c_surf[cat_idx] = exp(clamp(ui, +/-30))|.
\end{enumerate}

\section{What this transition did and did not change}

\textbf{Did change.}
\begin{itemize}
  \item Primary unknown $c_i \to u_i = \ln c_i$ on the dynamic species.
  \item Bulk Dirichlet BC type $c_i = c_{0,i} \to u_i = \ln c_{0,i}$;
        $c_{0,i}=0$ now requires a floor.
  \item Time stepping is backward Euler in $c = e^u$, so the time term is
        $(\,e^{u^{n+1}} - e^{u^n}\,)/\Delta t$, not $e^{u^{n+1}}(u^{n+1} - u^n)/\Delta t$.
  \item Concentrations are positive by construction, so the negative-$c$
        Newton iterates and the associated floor regularizations
        in the old solver are gone.
  \item Conditioning of Debye-layer regions is improved: the exponential
        in $c(x)$ becomes linear in $u(x)$.
  \item BV residual gains an algebraically rearranged log-rate option
        (\verb|bv_log_rate=True|) that lets the surface $u_i$ enter
        unclamped.
\end{itemize}

\textbf{Did not change.}
\begin{itemize}
  \item The Poisson PDE structure (only its source uses $e^{u_i}$).
  \item The BV reaction rate physics, $\Eeq$ values, $\alpha$, $k_0$,
        stoichiometry, or the meaning of $\eta$.
  \item The $\eta$-clip $\pm 50$ on $\eta_{\mathrm{scaled}}$, which is
        floating-point safety on $\exp(\alpha n_e \eta)$ and is structural,
        not a c-vs-u choice.
  \item The bulk $u$-clamp $\pm 30$, which is the necessary log-c analog
        of guarding $\exp(u)$ as a coefficient and \emph{cannot} be
        eliminated by the log-rate trick (different role: coefficient vs
        single deferred $\exp$).
  \item The Boltzmann counterion (\verb|add_boltzmann_counterion_residual|)
        is added orthogonally to the log-c transform.
\end{itemize}

\section{Where to read more}
\begin{itemize}
  \item \texttt{docs/PNP Equation Formulations.tex} --- original
        concentration formulation.
  \item \texttt{Forward/bv\_solver/forms\_logc.py} --- production
        implementation.
  \item \texttt{docs/clipping\_conventions.md} --- detailed accounting of
        the three clips/clamps and why each one stays.
  \item \texttt{docs/bv\_solver\_unified\_api.md} --- how to call the
        production stack via config flags.
\end{itemize}

\end{document}
